% ===== Джерела до розділу 6. Ключі з префіксом ch06: =====
% Позиції ядра (bib-core.tex) тут НЕ дублюються.

% --- Сигнатури та правила ---
\bibitem{ch06:yara-doc} Álvarez V. M. YARA Documentation: Writing YARA rules. Release~4.x. VirusTotal. URL: \url{https://yara.readthedocs.io/}.

\bibitem{ch06:clamav} Cisco-Talos. ClamAV: Open Source Antivirus Engine --- Signatures and Database Formats. URL: \url{https://docs.clamav.net/}.

\bibitem{ch06:christodorescu} Christodorescu M., Jha S., Seshia S. A., Song D., Bryant R. E. Semantics-Aware Malware Detection. \emph{2005 IEEE Symposium on Security and Privacy}. Oakland: IEEE, 2005. P.~32--46. DOI: 10.1109/SP.2005.20.

% --- Дизасемблювання та аналіз коду ---
\bibitem{ch06:eagle} Eagle C. The IDA Pro Book. 2nd ed. San Francisco: No Starch Press, 2011. 672~p.

\bibitem{ch06:andriesse} Andriesse D., Chen X., van der Veen V., Slowinska A., Bos H. An In-Depth Analysis of Disassembly on Full-Scale x86/x64 Binaries. \emph{Proceedings of the 25th USENIX Security Symposium}. USENIX Association, 2016. P.~583--600.

\bibitem{ch06:brumley} Brumley D., Jager I., Avgerinos T., Schwartz E. J. BAP: A Binary Analysis Platform. \emph{Computer Aided Verification (CAV 2011)}. LNCS Vol.~6806. Berlin: Springer, 2011. P.~463--469.

% --- Подібність та кластеризація ---
\bibitem{ch06:bindiff} Dullien T., Rolles R. Graph-based comparison of Executable Objects. \emph{Symposium sur la sécurité des technologies de l'information et des communications (SSTIC)}. 2005.

\bibitem{ch06:jang} Jang J., Brumley D., Venkataraman S. BitShred: Feature Hashing Malware for Scalable Triage and Semantic Analysis. \emph{Proceedings of the 18th ACM Conference on Computer and Communications Security (CCS~'11)}. New York: ACM, 2011. P.~309--320.

\bibitem{ch06:kolter} Kolter J. Z., Maloof M. A. Learning to Detect and Classify Malicious Executables in the Wild. \emph{Journal of Machine Learning Research}. 2006. Vol.~7. P.~2721--2744.

\bibitem{ch06:tlsh} Oliver J., Cheng C., Chen Y. TLSH --- A Locality Sensitive Hash. \emph{2013 Fourth Cybercrime and Trustworthy Computing Workshop}. Sydney: IEEE, 2013. P.~7--13.

% --- Українські джерела ---
\bibitem{ch06:hybrid} Метод захисту програмного забезпечення на основі гібридного аналізу коду. \emph{Кібербезпека: освіта, наука, техніка}. 2025. URL: \url{https://csecurity.kubg.edu.ua/index.php/journal/article/view/871}.

\bibitem{ch06:sast} Трофименко О. Г., Дика А. І., Лобода Ю. Г. Порівняльний аналіз методів тестування безпеки застосунків (SAST, DAST, IAST). \emph{Кібербезпека: освіта, наука, техніка}. 2023. Т.~4, №~20. URL: \url{https://csecurity.kubg.edu.ua/index.php/journal/article/view/467}.

\bibitem{ch06:ml-sast} Опірський І. Р., Мельничук М. Р. Застосування машинного навчання в статичному аналізі безпеки коду. \emph{Кібербезпека: освіта, наука, техніка}. 2026. URL: \url{https://csecurity.kubg.edu.ua/index.php/journal/article/view/1193}.
